
\begin{thebibliography}{999}
\raggedright %
\phantomsection
\addcontentsline{toc}{chapter}{参考文献}

\bibitem{bib1} C. Abbas and H. Hofer. Holomorphic curves and global questions in
contact geometry. To appear in Birkhäuser.

\bibitem{bib2} R. Abraham and J. Robbin. \emph{Transversal mappings and flows.}
W.A. Benjamin, Inc., 1967.

\bibitem{bib3} B. Aebischer, M. Borer, Ch. Leuenberger, M. Kälin, and
H.M. Reimann. \emph{Symplectic geometry, an introduction based on
the Seminar in Berne 1992.} Birkhäuser, 1994.

\bibitem{bib4} S. Alpern and V.S. Prasad. Combinatorial proofs of
the Conley-Zehnder-Franks theorem on
a fixed point for torus homeomorphisms. \emph{Advances in Mathematics},
99:238--247, 1993.

\bibitem{bib5} H. Amann. Lectures on some fixed point theorems.
Monografias de matemática IMPA, Rio de Janeiro.

\bibitem{bib6} A. Ambrosetti and P. Rabinowitz. Dual variational methods in
critical point theory and applications. \emph{J. Funct. Anal.},
14:349--381, 1973.

\bibitem{bib7} V.I. Arnold. Proof of a theorem of A.N. Kolmogorov on
the invariance of quasiperiodic motions under small perturbations of
the Hamiltonian. \emph{Uspeki Mat. Nauk}, 18:13--40, 1963.
Russian Math. Surveys 18 (1963), 9--36.

\bibitem{bib8} V.I. Arnold. Small denominators and problems of stability of
motions in classical and celestial mechanics. \emph{Uspeki Mat. Nauk},
18:91--196, 1963. Russian Math. Surveys 18 (1963), 85--193.

\bibitem{bib9} V.I. Arnold. Sur une propriété topologique des applications
globalement canoniques de la méchanique classique. \emph{C. R. Acad. Paris},
261:3719--3722, 1965.

\bibitem{bib10} V.I. Arnold. Proceedings of Symposia in Pure Mathematics.
Mathematical developments arising from Hilbert problems.
\emph{AMS, Providence}, 28:66, 1976.

\bibitem{bib11} V.I. Arnold. \emph{Mathematical methods of classical
mechanics.} Springer-Verlag, Berlin und New York, 1978.
In particular: Appendix 9.

\bibitem{bib12} V.I. Arnold and A. Avez. \emph{Ergodic problems of classical
mechanics.} W.A. Benjamin, Inc., 1968. In particular Appendix 26.

\bibitem{bib13} M. Audin and J. Lafontaine. Holomorphic curves in symplectic
geometry. Birkhäuser, 1994

\bibitem{bib14} V. Bangert. Geodätische Linien auf Riemannschen
Mannigfaltigkeiten. \emph{Jber. d. Dt. Math.-Verein.}, 87:39--66, 1985.

\bibitem{bib15} V. Bangert. On the existence of closed geodesics on two
spheres. \emph{Intern. Journal of Mathematics}, 4:1:1--10, 1993.

\bibitem{bib16} A. Banyaga. Sur la structure du groupe des difféomorphismes
qui préservent une forme symplectique. \emph{Comment. Math. Helvetici},
53:174--227, 1978.

\bibitem{bib17} W. Beckner. Sharp Sobolev inequalities on the sphere and
the Moser-Trudinger inequality. \emph{Annals of Mathematics},
138:213--242, 1993.

\bibitem{bib18} V. Benci. Closed geodesic for the Jacobi metric and periodic
solutions of prescribed energy of natural Hamiltonian systems.
\emph{Ann. Inst. H. Poincaré, Analyse non linéaire}, 1:401--412, 1984.

\bibitem{bib19} V. Benci, H. Hofer, and P. Rabinowitz. A remark on a-priori
bounds and existence for periodic solutions of Hamiltonian systems.
\emph{NATO ASI Series, Series C}, In: Periodic Solutions of Hamiltonian
Systems and Related Topics,(209):85--88, 1987.

\bibitem{bib20} D. Bennequin. Entrelacements et équations de Pfaff.
\emph{Astérisque}, 107--108:83--161, 1983.

\bibitem{bib21} H. Berestycki, J.M. Lasry, G. Mancini, and B. Ruf. Existence
of multiple periodic orbits on star-shaped Hamiltonian surfaces.
\emph{Comm. Pure Appl. Math.}, 38:253--289, 1985.

\bibitem{bib22} M. Bialy and L. Polterovich. Geodesics of Hofer's metric on
the group of Hamiltonian diffeomorphisms. To appear: Duke Math. Journal,
1993.

\bibitem{bib23} G.D. Birkhoff. Proof of Poincaré's geometric theorem.
\emph{Trans. Amer. Math. Soc.}, 14:14--22, 1913.

\bibitem{bib24} G.D. Birkhoff. The restricted problem of three bodies.
\emph{Rend. Circ. Mat. Palermo}, 39:265--334, 1915.

\bibitem{bib25} G.D. Birkhoff. \emph{Dynamical Systems.} Colloquium Publ.
A.M.S., Vol. 9, 1927.

\bibitem{bib26} G.D. Birkhoff. Stability and the equations of dynamics.
\emph{Am. J. Math.}, 49:1--38, 1927.

\bibitem{bib27} G.D. Birkhoff and D.C. Lewis. On the periodic motions near
a given periodic motion of a dynamical system. \emph{Annali di Math.},
12:117--133, 1933.

\bibitem{bib28} Ph. Blanchard and E. Brüning. \emph{Variational methods in
mathematical physics.} Springer, 1992.

\bibitem{bib29} S.V. Bolotin. Libration motions of natural dynamical systems.
\emph{Moscow University Bulletin}, 3, 1978.

\bibitem{bib30} Yu.G. Borisovich, V.G. Zvyagin, and Yu.I. Sapronov.
Non-linear Fredholm maps and the Leray-Schauder theory.
\emph{Russian Math. Surveys}, 32(4):1--54, 1977.

\bibitem{bib31} J. Bourgain. Approximation of solutions of the cubic
nonlinear Schrödinger equations by finite-dimensional equations and
nonsqueezing properties. \emph{IMRN}, 2:79--90, 1994.

\bibitem{bib32} H. Brezis and L. Nirenberg. Nonlinear functional analysis.
In preparation.

\bibitem{bib33} L.E.J. Brouwer. Beweis des ebenen Translationssatzes.
\emph{Math. Ann.}, 72:37--54, 1912.

\bibitem{bib34} M. Brown. A new proof of Brouwer's lemma on translation arcs.
\emph{Houston J. of Math.}, 10:35--41, 1984.

\bibitem{bib35} M. Brown and W.D. Neumann. Proof of the Poincaré-Birkhoff
fixed point theorem. \emph{Michigan Math. J.}, 24:21--31, 1977.

\bibitem{bib36} M. Brunella. On a theorem of Sikorav.
\emph{L'Enseignement Mathématique}, 37:83--87, 1991.

\bibitem{bib37} T. Carleman. Sur un problème d'unicité pour les systèmes
d'équations aux dérivées partielles à deux variables indépendantes.
\emph{Ark. Mat. Asts. Fys.}, 26B, No. 17:1--9, 1939.

\bibitem{bib38} P.H. Carter. An improvement of the Poincaré-Birkhoff theorem.
\emph{Trans. Amer. Math. Soc.}, 269:285--299, 1982.

\bibitem{bib39} K.-C. Chang. \emph{Infinite dimensional Morse theory and
multiple solution problems.} Birkhäuser, 1993.

\bibitem{bib40} S-Y. A. Chang. Extremal functions in a sharp form of Sobolev
inequality. \emph{ICM Berkeley}, 1:715--723, 1986.

\bibitem{bib41} M. Chaperon. Phases génératrices en géométrie symplectique.
In: D. Bennequin, Rencontre entre physiciens théoriciens et
mathématiciens en l'honneur de René Thom. Publication of the
Université de Strasbourg.

\bibitem{bib42} M. Chaperon. Quelques questions de géométrie symplectique.
\emph{Séminaire Bourbaki 1982--83, Astérisque}, 105--106:231--249, 1983.

\bibitem{bib43} M. Chaperon. Une idée du type ``géodésique brisée'' pour
les systèmes Hamiltoniens. \emph{C.R. Acad. Sci. Paris Sér. I Math.},
298:293--296, 1984. Série I, no. 13.

\bibitem{bib44} M. Chaperon. Lois de conservation et géométrie symplectique.
\emph{C.R. Acad. Sci. Paris}, 312:345--348, 1991. Série I.

\bibitem{bib45} M. Chaperon and B. D'Onofrio. Sur un théorème de Claude
Viterbo. \emph{C.R. Acad. Sci. Paris}, 306:597--600, 1988. Série I.

\bibitem{bib46} M. Chaperon and E. Zehnder. Quelques résultat en géométrie
symplectique. In \emph{Séminaire Sud-Rhodanien de Géométrie III: Autour
du Théorème de Poincaré-Birkhoff}, pages 51--121, 1984.
Travaux en cours. Hermann Paris.

\bibitem{bib47} Yu.V. Chekanov. Legandrova teorija morsa.
\emph{Uspekhi Mat. Naul.}, 42:4:129--141, 1987.

\bibitem{bib48} A. Chenciner. Sur un énoucé dissipatif du théorème
géométrique de Poincaré-Birkhoff. \emph{C.R. Acad. Sci. Paris},
294, I:243--245, 1982.

\bibitem{bib49} A. Chenciner. La dynamique on voisinage d'un point fixe
elliptique conservatif: de Poincaré et Birkhoff à Aubry et Mather.
\emph{Séminaire Bourbaki Astérisque}, 121/122:147--170, 1984.

\bibitem{bib50} A. Floer, H. Hofer and K. Wysocki. Applications of symplectic
homology I, to appear in \emph{Math. Zeit.}

\bibitem{bib51} K. Cieliebak, A. Floer, and H. Hofer. Symplectic homology II:
A general construction. To appear in Math. Zeit.

\bibitem{bib52} K. Cieliebak, A. Floer, H. Hofer, and K. Wysocki.
Applications of symplectic homology II. Preprint E.T.H. Zürich (1994).

\bibitem{bib53} F. Clarke. A classical variational principle for periodic
Hamiltonian trajectories. \emph{Proc. Amer. Math. Soc.}, 76:186--188, 1979.

\bibitem{bib54} C. Conley. \emph{Isolated invariant sets and Morse index.}
CBMS Reg. Conf. Series in Math. 38. AMS, 1978.

\bibitem{bib55} C. Conley and E. Zehnder. The Birkhoff-Lewis fixed point
theorem and a conjecture of V. I.Arnold. \emph{Invent. Math.},
73:33--49, 1983.

\bibitem{bib56} C. Conley and E. Zehnder. \emph{An index theory for periodic
solutions of a Hamiltonian system}, volume 1007 of \emph{Geometric
Dynamics, Lecture Notes in Mathematics.} Springer, 1983. 132--145.

\bibitem{bib57} C. Conley and E. Zehnder. Morse type index theory for flows
and periodic solutions of Hamiltonian equations.
\emph{Comm. Pure Appl. Math.}, 37:207--253, 1984.

\bibitem{bib58} C.B. Croke and A. Weinstein. Closed curves on convex
hypersurfaces and periods of nonlinear oscillations.
\emph{Invent. Math.}, 64:199--202, 1981.

\bibitem{bib59} B. Dacorogna and J. Moser. On a partial differential equation
involving the Jacobian determinant. \emph{Ann. Inst. Henri Poincaré,
Analyse non linéaire}, 7:1--26, 1990.

\bibitem{bib60} K. Deimling. \emph{Nonlinear functional analysis.}
Springer-Verlag, Berlin Heidelberg, 1985.

\bibitem{bib61} W.-Y. Ding. A generalization of the Poincaré-Birkhoff theorem.
\emph{Proc. AMS}, 88:341--346, 1983.

\bibitem{bib62} A. Dold. \emph{Lectures on Algebraic Topology.} Grundlagen
der mathematischen Wissenschaften 200, Springer 1980. second edition.

\bibitem{bib63} A. Douglis and L. Nirenberg. Interior estimates for elliptic
systems of partial differential equations. \emph{Comm. Pure Appl. Math.},
8:503--538, 1955.

\bibitem{bib64} J.J. Duistermaat. On global action-angle coordinates.
\emph{Commun. on Pure and Appl. Math.}, 33:687--706, 1980.

\bibitem{bib65} I. Ekeland. Une théorie de Morse pour les systèmes
hamiltoniens convexes. \emph{Ann. Inst. Henri Poincaré, Analyse non
linéaire}, 1:19--78, 1984.

\bibitem{bib66} I. Ekeland. \emph{Convexity methods in Hamiltonian
mechanics.} Springer, 1990. Ergebnisse der Mathematik und ihrer
Grenzgebiete, 3. Folge.

\bibitem{bib67} I. Ekeland and H. Hofer. Convex Hamiltonian energy surfaces
and their closed trajectories. \emph{Comm. Math. Physics}, 113:419--467,
1987.

\bibitem{bib68} I. Ekeland and H. Hofer. Symplectic topology and Hamiltonian
dynamics. \emph{Math. Zeit.}, 200:355--378, 1990.

\bibitem{bib69} I. Ekeland and H. Hofer. Symplectic topology and Hamiltonian
dynamics II. \emph{Math. Zeit.}, 203:553--567, 1990.

\bibitem{bib70} I. Ekeland and J.M. Lasry. On the number of periodic
trajectories for a Hamiltonian flow on
a convex energy surface. \emph{Ann. Math}, 112:283--319, 1980.

\bibitem{bib71} Y. Eliashberg. A theorem on the structure of wave fronts and
its application in symplectic topology. \emph{Funct. Anal. and its Appl.},
21:227--232, 1987.

\bibitem{bib72} Y. Eliashberg. Topological characterization of Stein
manifolds of dimension $> 2$. \emph{Intern. J. Math.}, 1:29--46, 1990.

\bibitem{bib73} Y. Eliashberg. Filling by holomorphic discs and its
applications. \emph{London Math. Society Lecture Notes}, pages 45--67,
1991. Series 151.

\bibitem{bib74} Y. Eliashberg. Contact 3-manifolds, twenty year since
J. Martinet's work. \emph{Ann. Inst. Fourier}, 42:165--192, 1992.

\bibitem{bib75} Y. Eliashberg and H. Hofer. Unseen symplectic boundaries.
In \emph{Proceedings of a Conference in honour of E. Calabi}, 1991.

\bibitem{bib76} Y. Eliashberg and L. Polterovich. Bi-invariant metrics on
the group of Hamiltonian diffeomorphisms. International Journal of
Mathematics Vol. 4, No. 5 (1993) 727--738.

\bibitem{bib77} Y.M. Eliashberg. Estimates on the number of fixed points of
area preserving transformations. Preprint Syktyvkar 1978 (Russian).

\bibitem{bib78} Y.M. Eliashberg. Rigidity of symplectic and contact
structures. Preprint, 1981.

\bibitem{bib79} H. Eliasson. Geometry of manifolds of maps.
\emph{J. Diff. Geom.}, 1:169--194, 1967.

\bibitem{bib80} A. Fahti. An orbit closing proof of Brouwer's lemma on
translation arcs. \emph{Enseignement Math.}, 33:315--322, 1987.

\bibitem{bib81} A.I. Fet and L.A. Lusternik. Variational problems on closed
manifolds. \emph{Dokl. Akad. Nauk. S.S.S.R.}, 81:17--18, 1951.

\bibitem{bib82} A. Floer. Proof of the Arnold conjecture for surfaces and
generalizations to certain Kähler manifolds. \emph{Duke Math. J.},
53:1--32, 1986.

\bibitem{bib83} A. Floer. An instanton invariant for 3-manifolds.
\emph{Comm. Math. Phys.}, 118:215--240, 1988.

\bibitem{bib84} A. Floer. Morse theory for Lagrangian intersection theory.
\emph{J. Diff. Geom.}, 18:513--517, 1988.

\bibitem{bib85} A. Floer. The unregularised gradient flow of the symplectic
action. \emph{Comm. Pure Appl. Math.}, 41:775--813, 1988.

\bibitem{bib86} A. Floer. Cuplength estimates on Lagrangian intersections.
\emph{Comm. Pure Appl. Math.}, 42:335--356, 1989.

\bibitem{bib87} A. Floer, Hofer H., and D. Salamon. Transversality in
the elliptic Morse theory for the symplectic action.
Preprint E.T.H. Zürich 1994.

\bibitem{bib88} A. Floer, Hofer H., and C. Viterbo. The Weinstein conjecture
in $P \times \mathbb{C}^l$. \emph{Math. Zeit.}, 203:355--378, 1989.

\bibitem{bib89} A. Floer and H. Hofer. Coherent orientations for periodic
orbit problems in symplectic geometry. \emph{Math. Zeit.}, 212:13--38,
1993.

\bibitem{bib90} A. Floer and H. Hofer. Symplectic homology I: Open sets in
$\mathbb{C}^n$. \emph{Math. Zeit.}, 215:37--88, 1994.

\bibitem{bib91} M. Flucher. Fixed points of measure preserving torus
homeomorphisms. \emph{manuscripta math.}, 68:271--293, 1990.

\bibitem{bib92} O. Forster. \emph{Riemannsche Flächen.} Heidelberger
Taschenbücher. Springer, 1977.

\bibitem{bib93} B. Fortune. A symplectic fixed point theorem for $CP^n$.
\emph{Invent. Math.}, 81:22--46, 1985.

\bibitem{bib94} B. Fortune and A. Weinstein. A symplectic fixed point theorem
for complex projective spaces. \emph{Bulletin AMS}, 12:128--130, 1985.

\bibitem{bib95} J. Franks. Periodic points and rotation numbers for area
preserving diffeomorphisms of the plane. Preprint, Northwestern University.

\bibitem{bib96} J. Franks. Generalizations of the Poincaré-Birkhoff theorem.
\emph{Ann. of Math.}, 128:139--151, 1988.

\bibitem{bib97} J. Franks. Recurrence and fixed points of surface
homeomorphisms. \emph{Ergod. Th. and Dynam. Sys.}, 8:99--107, 1988.
Charles Conley Memorial Issue.

\bibitem{bib98} J. Franks. \emph{A variation of the Poincaré-Birkhoff
theorem.} Amer. Math. Soc., Providence, 1988. Contemp. Math., 81.

\bibitem{bib99} J. Franks. Periodic points and rotation numbers for area
preserving diffeomorphisms of the plane. \emph{IHES-Publ.Math.},
71:105--120, 1990.

\bibitem{bib100} J. Franks. Geodesics on $S^2$ and periodic points of
annulus homeomorphisms. \emph{Invent. Math.}, 108:403--418, 1992.

\bibitem{bib101} J. Franks. A new proof of the Brouwer plane translation
theorem. \emph{Ergod. Th. and Dynam. Sys.}, 12:217--226, 1992.

\bibitem{bib102} K. Fukaya. Floer homology for oriented 3-manifolds.
\emph{Advanced Studies in Pure Mathematics}, 20:1--92, 1992.

\bibitem{bib103} J.M. Gambaudo and P. Le Calvez. Infinié d'orbites
périodiques pour les difféomorphismes conservatifs de l'anneau.
Preprint, 1991.

\bibitem{bib104} C. Genecand. Transversal homoclinic orbits near elliptic
fixed points of area preserving diffeomorphisms of the plane.
\emph{Dynamics Reported, New Series}, 2:1--30, 1993.

\bibitem{bib105} A.B. Givental. Global properties of Maslov index and Morse
theory. \emph{Funct. Anal. and its Appl.}, 22:141--142, 1988.

\bibitem{bib106} H. Gluck and W. Ziller. \emph{Existence of periodic motions
of conservative systems.} Princeton University Press, 1983. In: Seminar
on minimal submanifolds, E. Bombieri, ed.

\bibitem{bib107} M. Gromov. Pseudoholomorphic curves in symplectic manifolds.
\emph{Invent. Math.}, 82:307--347, 1985.

\bibitem{bib108} M. Gromov. \emph{Partial differential relations.}
Ergebnisse der Mathematik. Springer, 1986.

\bibitem{bib109} L. Guillon. Le théorème de translation de Brouwer: une
démonstration simplifiée menant
à une nouvelle preuve du théorème de Poincaré-Birkhoff. 1991.

\bibitem{bib110} M.-R. Herman. Differentiabilite optimale et contre-exemples
à la fermeture en topologie $C^\infty$ des orbites
recurrentes de flots Hamiltoniens. \emph{Comptes Rendus de l'Académie des
Sciences. Serie-I.-Mathématique}, 313:49--51, 1991. No. 1.

\bibitem{bib111} M.-R. Herman. Exemples de flots Hamiltoniens dont aucune
perturbation en topologie $C^\infty$ n'a d'orbites periodiques sur un
ouvert de surfaces d'energies. \emph{Comptes Rendus de l'Académie des
Sciences. Serie-I.-Mathématique}, 312:989--994, 1991. No. 13.

\bibitem{bib112} D. Hilbert. Über das Dirichlet'sche Prinzip.
\emph{Jhrb. der Deutschen Math. Vereinigung}, VIII, Erstes Heft:184--188,
1900.

\bibitem{bib113} M.W. Hirsch. \emph{Differential topology.} Graduate Texts
in Math. 33. Springer-Verlag, 1976.

\bibitem{bib114} H. Hofer. Lagrangian embeddings and critical point theory.
\emph{Ann. Inst. H. Poincaré, Anal. Non Linéaire}, 2:407--462, 1985.

\bibitem{bib115} H. Hofer. Ljusternik-Schnirelman theory for Lagrangian
intersections. \emph{Ann. Inst. H. Poincaré}, 5(5):465--499, 1988.

\bibitem{bib116} H. Hofer. On the topological properties of symplectic maps.
\emph{Proceedings of the Royal Society of Edinburgh}, 115 A:25--38, 1990.

\bibitem{bib117} H. Hofer. Estimates for the energy of a symplectic map.
\emph{Comm. Math. Helv.}, 68:48--72, 1993.

\bibitem{bib118} H. Hofer. Pseudoholomorphic curves in symplectisations
with application to the Weinstein conjecture in dimension three.
\emph{Invent. Math.}, 114:515--563, 1993.

\bibitem{bib119} H. Hofer and D. Salamon. Floer homology and Novikov rings.
To appear A. Floer Memorial Volume.

\bibitem{bib120} H. Hofer, C. Taubes, A. Weinstein, and E. Zehnder, editors.
\emph{Floer Memorial Volume.} To appear, Birkhäuser, 1994.

\bibitem{bib121} H. Hofer and C. Viterbo. The Weinstein conjecture in
cotangent bundles and related results. \emph{Annali Sc. Norm. Sup. Pisa},
15:411--445, 1988. Serie IV, Fasc III.

\bibitem{bib122} H. Hofer and C. Viterbo. The Weinstein conjecture in the
presence of holomorphic spheres. \emph{Comm. Pure Appl. Math.},
45(5):583--622, 1992.

\bibitem{bib123} H. Hofer and E. Zehnder. \emph{A new capacity for symplectic
manifolds}, pages 405--428. Analysis et cetera. Academic press, 1990.
Edit: P. Rabinowitz and E. Zehnder.

\bibitem{bib124} L. Hörmander. \emph{The analysis of linear partial
differential operators}, volume I--IV. Springer, Grundlehren Edition,
1983--1985.

\bibitem{bib125} M.-Y. Jiang. Hofer-Zehnder symplectic capacity for two
dimensional manifolds. \emph{Proc. Royal Soc. Edinb.}, 123A:945--950, 1993.

\bibitem{bib126} R. Jost. Winkel- und Wirkungsvariable für allgemeine
mechanische Systeme. \emph{Helvetica Physica Acta}, 41:965--968, 1968.

\bibitem{bib127} W. Klingenberg. \emph{Lectures on closed geodesics.}
Springer, Grundlehren Math. Wiss. 230, 1978.

\bibitem{bib128} A.N. Kolmogorov. \emph{The general theory of dynamical
systems and classical mechanics.} Proc. Int. Congress of Math.
Amsterdam 1954. W.A. Benjamin, Inc., 1967. see Appendix D in R. Abraham:
Foundations of mechanics.

\bibitem{bib129} V.V. Koslov. Calculus of variations in the large and
classical mechanics. \emph{Russian Math. Survey}, 40:2:37--71, 1985.

\bibitem{bib130} S. Kuksin. Infinite-dimensional capacities and a squeezing
theorem for Hamiltonian PDE's. Preprint F.I.M., ETH Zürich (1993).

\bibitem{bib131} A. Künzle. The least characteristic action as symplectic
capacity. F.I.M. report, ETH Zürich (1991).

\bibitem{bib132} K. Kuperberg. A smooth counter example to the Seifert
conjecture. To appear in Annals of Mathematics.

\bibitem{bib133} F. Lalonde. Energy and capacities in symplectic topology.
Preprint, 1994.

\bibitem{bib134} F. Lalonde and D. McDuff. The geometry of symplectic
energy. Preprint, 1993.

\bibitem{bib135} F. Lalonde and D. McDuff. Hofer's $L^\infty$-geometry:
energy and stability of Hamiltonian flows I. Preprint, 1994.

\bibitem{bib136} F. Laudenbach and J.C. Sikorav. Persistence d'intersections
avec la section nulle au cours d'une isotopie Hamiltonienne dans un
fibre cotangent. \emph{Invent. Math.}, 82:349--357, 1985.

\bibitem{bib137} H.V. Lê and Kaoru Ono. Symplectic fixed points, Calabi
invariant and Novikov homology. Max-Planck-Institut Bonn MPI/93--27.

\bibitem{bib138} P. Le Calvez. Sur les difféomorphismes de l'anneau et du
tore. Preprint, Université Paris-Sud.

\bibitem{bib139} P. Le Calvez. Une généralisation du théorème de
Conley-Zehnder aux homéomorphismes du tore de dimension deux.
Prépublications Mathématiques de l'Université Paris-Nord, 93--02, 1993.

\bibitem{bib140} P. Le Calvez. Phases génératrices discrètes de
difféomorphismes du tore ou de l'anneau. Preprint, Unversité Paris-Sud,
1991.

\bibitem{bib141} P. Le Calvez. Un théorème de translations pour les
difféomorphismes du tore ayant des points fixes. Prépublications
Université de Paris-Sud, mathématiques, 92--85 (1992), 1991.

\bibitem{bib142} E. Lima. Orientability of smooth hypersurfaces and the
Jordan-Brouwer separation theorem. \emph{Expo. Math.}, 5:283--286, 1987.

\bibitem{bib143} L. Lusternik and L. Schnirelman. \emph{Méthodes
topologiques dans les problèmes variationnels.} Herman, Paris, 1934.

\bibitem{bib144} R. Lutz. Structures de contact sur les fibrés principaux
en circles de dimension 3. \emph{Ann. Inst. Fourier}, 3:1--15, 1977.

\bibitem{bib145} A.M. Lyapunov. Problème général de la stabilité du
movement. \emph{Ann. Fac. Sci. Toulouse}, 2:203--474, 1907.

\bibitem{bib146} R. Ma. A remark on Hofer-Zehnder symplectic capacity in
$M \times \mathbb{R}^{2n}$, 1992. Preprint, Tsinghua University.

\bibitem{bib147} J. Martinet. Formes de contact sur les variétés de
dimension 3. \emph{Lecture Notes in Math.}, 209:142--163, 1971.

\bibitem{bib148} W.S. Massey. \emph{Algebraic topology: an introduction.}
Graduate Texts in Mathematics 56. Springer Verlag, 1967.

\bibitem{bib149} W.S. Massey. \emph{Homology and cohomology theory.}
Marcel Dekker, Incl. New York and Basel, 1978.

\bibitem{bib150} W.S. Massey. \emph{Singular homology theory.} Graduate
Texts in Mathematics 70. Springer Verlag, 1980.

\bibitem{bib151} J. Mather. Existence of quasiperiodic orbits for twist
homeomorphisms of the annulus. \emph{Topology}, 21:457--467, 1982.

\bibitem{bib152} J. Mawhin and M. Willem. \emph{Critical point theory and
Hamiltonian systems.} Springer, 1989.

\bibitem{bib153} D. McDuff. Elliptic methods in symplectic geometry.
\emph{Bulletin AMS}, 23(2):311--358, 1990.

\bibitem{bib154} D. McDuff. Symplectic manifolds with contact type
boundaries. \emph{Invent. Math.}, 103:651--671, 1991.

\bibitem{bib155} D. McDuff and L. Polterovich. Symplectic packings and
algebraic geometry. Preprint, 1993.

\bibitem{bib156} D. McDuff and D. Salamon. Introduction to symplectic
topology. To appear.

\bibitem{bib157} J. Milnor. \emph{Topology from the differentiable
viewpoint.} The University Press of Virginia, Charlottesville, 1972.

\bibitem{bib158} A.S. Mishchenko and A.T. Fomenko. Generalized Liouville
method of integration of Hamiltonian systems. \emph{Funct. Anal. and
its Appl.}, 12:113--121, 1978.

\bibitem{bib159} J. Moser. On invariant curves of area preserving mappings
of the annulus. \emph{Nachr. Akad. Wiss. Göttingen, Math. Phys. Kl.},
pages 1--20, 1962.

\bibitem{bib160} J. Moser. On the volume elements on a manifold.
\emph{Trans. Amer. Math. Soc.}, 120:286--294, 1965.

\bibitem{bib161} J. Moser. A rapidly convergent iteration method and
non-linear partial differential equations I and II. \emph{Ann. Scuola
Norm. Sup. Pisa}, 20:265--315, 499--535, 1966.

\bibitem{bib162} J. Moser. \emph{Stable and random motions in dynamical
systems.} Princeton University Press, 1973. In: Annals of Math. Studies,
Vol. 77.

\bibitem{bib163} J. Moser. Periodic orbits near an equilibrium and a
theorem of A. Weinstein. \emph{Comm. Pure Appl. Math.}, 29:727--747, 1976.

\bibitem{bib164} J. Moser. Proof of a generalized form of a fixed point
theorem due to G.D. Birkhoff. \emph{Geometry and Topology, Lecture Notes
in Mathematics}, 597:464--494, 1977.

\bibitem{bib165} J. Moser. A fixed point theorem in symplectic geometry.
\emph{Acta Mathematica}, 141:17--34, 1978.

\bibitem{bib166} J. Moser. Integrable Hamiltonian systems and spectral
theory. In \emph{Lezioni Fermiane, Accademia nazionale dei Lincai
Scuola Normale Superiore, Pisa}, 1981.

\bibitem{bib167} J. Moser and E. Zehnder. \emph{Notes on Dynamical
systems.} Courant Institute, 1980.

\bibitem{bib168} N. Nekhoroshev. Action-angle variables and their
generalization. \emph{Trans. Mosc. Math. Soc.}, 26:181--198, 1972.

\bibitem{bib169} N. Nikishin. Fixed points of diffeomorphisms on the two
sphere that preserve area. \emph{Funkts. Anal. i Prilozhen}, 8:84--85,
1974.

\bibitem{bib170} S.P. Novikov. Multivalued functions and functionals ---
an analogue of the Morse theory. \emph{Soviet Math. Dokl.}, 24:222--225,
1981.

\bibitem{bib171} D. Offin. A class of periodic orbits in classical
mechanics. \emph{J. Diff. Eq.}, 66:90--116, 1987.

\bibitem{bib172} K. Ono. On the Arnold conjecture for weakly monotone
symplectic manifolds. Preprint Max-Planck-Institut für Mathematik
MPI/93--4.

\bibitem{bib173} R.S. Palais. Critical point theory and the minimax
principle. In: Global Analysis. \emph{Proc. AMS}, XV:185--212, 1970.

\bibitem{bib174} R.S. Palais and Terng C.-L. \emph{Critical point theory and
submanifold geometry,} volume 1353. Springer Lecture Notes in
Mathematics, 1988.

\bibitem{bib175} R.S. Palais and S. Smale. A generalized Morse theory.
\emph{Bull. AMS}, 70:165--171, 1964.

\bibitem{bib176} H. Poincaré. \emph{Methodes nouvelle de la méchanique
céleste,} volume 3. Gauthier-Villars, Paris, 1899.

\bibitem{bib177} H. Poincaré. Sur un théorème de géométrie.
\emph{Rend. Circolo Math. Palermo}, 33:375--407, 1912.

\bibitem{bib178} J. Pöschel. Integrability of Hamiltonian systems on
Cantor sets. \emph{Comm. Pure Appl. Math.}, 35:653--696, 1982.

\bibitem{bib179} C.C. Pugh and R.C. Robinson. The $C^1$-closing lemma,
including Hamiltonians. \emph{Erg. Th. and Dynam. Sys.}, 3:261--313, 1983.

\bibitem{bib180} P. Rabinowitz. Periodic solutions of Hamiltonian systems.
\emph{Comm. Pure Appl. Math.}, 31:157--184, 1978.

\bibitem{bib181} P. Rabinowitz. Periodic solutions of Hamiltonian systems
on a prescribed energy surface. \emph{J. Diff. Equ.}, 33:336--352, 1979.

\bibitem{bib182} P. Rabinowitz. \emph{Minimax methods in critical point
theory with applications to differential equations.} CBMS, Regional
Conf. Ser. in Math. 65. AMS, 1986.

\bibitem{bib183} C. Robinson. A global approximation theorem for
Hamiltonian systems. \emph{Proc. Symposium in Pure Math.}, XIV:233--244,
1970. Global Analysis, AMS.

\bibitem{bib184} K.P. Rybakowski. On the homotopy index for
infinite-dimensional semiflows. \emph{Trans. Amer. Math. Soc.},
269:351--382, 1982.

\bibitem{bib185} K.P. Rybakowski. \emph{The homotopy index and partial
differential equations.} Springer, 1987.

\bibitem{bib186} K.P. Rybakowski and E. Zehnder. A Morse equation in
Conley's index theory for semiflows on metric spaces. \emph{Ergod. Th.
and Dynam. Sys.}, 5:123--143, 1985.

\bibitem{bib187} D. Salamon. Connected simple system and the Conley
index of isolated invariant sets. \emph{Trans. Amer. Math. Soc.},
291:1--41, 1985.

\bibitem{bib188} D. Salamon. Morse theory, the Conley index and Floer
homology. \emph{Bull. London Math. Soc.}, 22:113--140, 1990.

\bibitem{bib189} D. Salamon and E. Zehnder. KAM theory in configuration
space. \emph{Comm. Math. Helvetici}, 64:84--132, 1989.

\bibitem{bib190} D. Salamon and E. Zehnder. Morse theory for periodic
solutions of Hamiltonian systems and the Maslov index.
\emph{Comm. Pure Appl. Math.}, 45:1303--1360, 1992.

\bibitem{bib191} M. Schwarz. \emph{Morse homology.} Progress in
Mathematics, 111. Birkhäuser, 1993.

\bibitem{bib192} P.A. Schweitzer. Counter examples to the Seifert
conjecture and opening closed leaves of foliations. \emph{Ann. Math.},
100:386--400, 1974.

\bibitem{bib193} H. Seifert. Periodische Bewegungen mechanischer Systeme.
\emph{Math. Zeit.}, 51:197--216, 1948.

\bibitem{bib194} K.F. Siburg. Symplectic capacities in two dimensions.
\emph{manuscripta math.}, 78:149--163, 1993.

\bibitem{bib195} C.L. Siegel. über die Existenz einer Normalform
analytischer Hamiltonischer Differentialgleichungen in der Nähe einer
Gleichgewichtslösung. \emph{Math. Ann.}, 128:144--170, 1954.

\bibitem{bib196} C.L. Siegel and J. Moser. \emph{Lectures on Celestial
Mechanics.} Springer Grundlehren, Band 187, 1971.

\bibitem{bib197} J.-C. Sikorav. Problèmes d'intersections et de points
fixes en géométrie Hamiltonienne. \emph{Comm. Math. Helv.}, 62:62--73,
1987.

\bibitem{bib198} J.-C. Sikorav. Systèmes Hamiltoniens et topologie
symplectique. Dipartimento di Matematica dell' Università di Pisa, 1990.
ETS, EDITRICE PISA.

\bibitem{bib199} J.-C. Sikorav. Homologie associée à une fonctionelle
(d'après A. Floer). \emph{Séminaire Bourbaki Exp. 730--744 Astérisque},
201--202--203:115--143, 1991.

\bibitem{bib200} C.P. Simon. A bound for the fixed point index
of an area preserving map with applications in mechanics.
\emph{Invent. Math.}, 26:187--200, 1974.

\bibitem{bib201} S. Smale. An infinite dimensional version of Sard's
theorem. \emph{Amer. J. Math.}, 87:861--866, 1965.

\bibitem{bib202} E. Spanier. Cohomology theory for general spaces.
\emph{Ann. of Math. (2)}, 49:407--427, 1948.

\bibitem{bib203} E. Spanier. \emph{Algebraic topology.} Springer Verlag,
New York, 1966.

\bibitem{bib204} E. Spanier. Tautness for Alexander-Spanier cohomology.
\emph{Pacific Journal of Mathematics}, (75) 2:561--563, 1978.

\bibitem{bib205} E.M. Stein. \emph{Singular integrals and differentiability
properties of functions.} Princeton University Press, 1970.

\bibitem{bib206} A. Stephen. On homeomorphisms of the plane which have no
fixed points. \emph{Abh. Math. Seminar Univ. Hamburg}, 30:61--74, 1967.

\bibitem{bib207} M. Struwe. Existence of periodic solutions of Hamiltonian
systems on almost every energy surface. \emph{Bol. Soc. Bras. Mat.},
20:49--58, 1990.

\bibitem{bib208} M. Struwe. \emph{Variational methods.} Springer, 1990.

\bibitem{bib209} R. Switzer. \emph{Algebraic topology --- homotopy and
homology.} Grundlehren, 212. Springer Verlag, Berlin, 1975.

\bibitem{bib210} C.H. Taubes. Self-dual Yang-Mills connections on
non-self-dual 4-manifolds. \emph{J. Diff. Geom.}, 17:139--170, 1982.

\bibitem{bib211} C.H. Taubes. A framework for Morse theory for the
Yang-Mills functionals. Preprint Harvard University, 1986.

\bibitem{bib212} L. Traynor. Symplectic homology via generating functions.
Preprint Stanford University, 1993.

\bibitem{bib213} L. Traynor. Symplectic packing constructions. Preprint,
1993.

\bibitem{bib214} I. Ustilovsky. Conjugate points on geodesics of Hofer's
metric. Preprint, 1993.

\bibitem{bib215} I.N. Vekua. \emph{Verallgemeinerte analytische
Funktionen.} Akademie-Verlag, Berlin, 1963.

\bibitem{bib216} J. Vick. \emph{Homology theory, an introduction to
algebraic topology.} Academic Press, 1973.

\bibitem{bib217} C. Viterbo. A proof of the Weinstein conjecture in
$\mathbb{R}^{2n}$. \emph{Ann. Inst. H. Poincaré, Anal. non linéaire},
4:337--357, 1987.

\bibitem{bib218} C. Viterbo. Symplectic topology as the geometry of
generating functions. \emph{Math. Annalen}, 292:685--710, 1992.

\bibitem{bib219} F. Warner. \emph{Foundations of differentiable manifolds and
Lie groups.} Scott, Foresman and Company, 1971.

\bibitem{bib220} K. Weierstrass. \emph{Mathematische Werke.} 1858. Berlin,
Band I: 233--246, Band II: 19--44, Nachtrag: 139--148.

\bibitem{bib221} A. Weinstein. Contact surgery and symplectic handlebodies.
\emph{Hokkaido Math. Journal} 20 (1990) 241--251.

\bibitem{bib222} A. Weinstein. Lagrangian submanifolds and Hamiltonian
systems. \emph{Ann. Math.}, 98:377--410, 1973.

\bibitem{bib223} A. Weinstein. Normal modes for nonlinear Hamiltonian
systems. \emph{Invent. Math.}, 20:47--57, 1973.

\bibitem{bib224} A. Weinstein. \emph{Lectures on symplectic manifolds},
volume 29 of \emph{CBMS, Reg. Conf. Series in Math.} AMS, 1977.

\bibitem{bib225} A. Weinstein. Periodic orbits for convex Hamiltonian
systems. \emph{Ann. Math.}, 108:507--518, 1978.

\bibitem{bib226} A. Weinstein. On the hypothesis of Rabinowitz's periodic
orbit theorems. \emph{J. Diff. Equ.}, 33:353--358, 1979.

\bibitem{bib227} A. Weinstein. On extending the Conley-Zehnder fixed point
theorem to other manifolds. \emph{Proc. Sympos. Pure Math.}, 45,
part 2:541--544, 1986. AMS.

\bibitem{bib228} F.W. Wilson. On the minimal sets of nonsingular vector
fields. \emph{Ann. Math.}, 84:529--536, 1966.

\bibitem{bib229} J.C. Yoccoz. Travaux de Herman sur les tores invariants.
\emph{Bourbaki Exp. 745--579, Astérisque}, 206:311--345, 1992.

\bibitem{bib230} E. Zehnder. Homoclinic points near elliptic fixed points.
\emph{Comm. Pure Appl. Math.}, 26:131--182, 1973.

\bibitem{bib231} E. Zehnder. Remarks on periodic solutions on hypersurfaces.
\emph{NATO ASI Series, Series C}, In: Periodic Solutions of Hamiltonian
Systems and Related Topics,(209):267--279, 1987.

\bibitem{bib232} E. Zeidler. \emph{Nonlinear functional analysis and its
applications I.} Springer-Verlag, New York Inc., 1986.

\end{thebibliography}
